\chapter*{Remerciements}
\addcontentsline{toc}{chapter}{Remerciements}
\markboth{REMERCIEMENTS}{REMERCIEMENTS}

La réalisation de ce projet de fin d’études et sa soutenance marque la fin du cycle Licence respectivement d’ingénieurs et inspecteurs des postes et télécommunications au sein de l’École Nationale Supérieure des Postes, des Télécommunications et des Technologies de l’Information et de la Communication de Yaoundé (SUP’PTIC). Différentes personnes ont contribué à la réalisation de ce projet. Nous tenons à remercier de façon particulière :

\begin{itemize}
    \item \textbf{M. FELIX WATCHING}, Directeur de l’école National Supérieur des Postes, télécommunications et des TIC, pour le cadre adéquat mis à notre disposition tout au long de cette formation ;
    
    \item \textbf{M. SASSA GIQUEL THERANCE}, notre encadreur académique, pour avoir accepté de nous encadrer, pour son accompagnement rigoureux, ses orientations méthodologiques et ses observations toujours pertinentes qui ont su guider notre réflexion tout au long de la réalisation de ce travail ;
    
    \item \textbf{M. NUEBISSI SIMO LANDRY}, notre encadreur professionnel, pour sa disponibilité, ses conseils techniques et sa capacité à traduire nos idées en réalités concrètes ;
    
    \item \textbf{Au corps enseignant et administratif de SUP’PTIC Yaoundé}, pour la qualité des enseignements et le cadre d’étude propice qui a été mis à notre disposition durant cette formation ;
    
    \item \textbf{Les membres de nos familles}, pour leur amour indéfectible et leur confiance qui ont été le socle de notre persévérance tout au long de ce parcours exigeant ;
    
    \item \textbf{Tous nos camarades de promotion} avec qui nous avons traversé de multiples épreuves ;
    
    \item \textbf{Tous ceux} qui de près ou de loin ont contribué à l’élaboration de ce travail et dont les noms ne figurent malheureusement pas sur ces pages. Qu’ils trouvent ici le témoignage de notre estime et de notre gratitude.
\end{itemize}